\SourceLeaf{401}%
oder
\[
\alpha=1,\quad\beta=0
\]
\par
Der Nenner in den Ausdrücken für \(\xi,\eta\) ist der doppelte Flächeninhalt des Dreiecks.

\par\medskip\Needspace{5\baselineskip}
\pdfbookmark[2]{II.}{carnot-zusatz-ii}
\centerline{\textbf{II.}}\nobreak\smallskip
Dass sich die Senkrechten von den Eckpunkten eines Dreiecks auf die gegenüberliegenden Seiten in einem Punkt schneiden, lässt sich sehr einfach zeigen.

\begingroup\rightskip=0pt plus 1em
Das gegebene Dreieck sei \(BDF\), und die genannten Senkrechten seien \(\overline{BI},\overline{DG},\overline{FH}\).\par
\endgroup

Man ziehe durch jeden Eckpunkt des Dreiecks eine Parallele zur gegenüberliegenden Seite. Diese Geraden schneiden sich in den Punkten \(A,C,E\). Folglich steht \(\overline{FH}\) auch auf \(\overline{AE}\) senkrecht, \(\overline{GD}\) auf \(\overline{CE}\) und \(\overline{BI}\) auf \(\overline{AC}\). Außerdem gilt
\[
\begin{aligned}\overline{AB}&=\overline{BC}\\\overline{ED}&=\overline{DC}\\\overline{AF}&=\overline{FE}\end{aligned}
\]
Beschreibt man nun einen Kreis um das Dreieck \(ACE\), so liegt sein Mittelpunkt sowohl auf \(\overline{BI}\) als auch auf \(\overline{DG}\) und auf \(\overline{FH}\). Diese drei Geraden müssen sich also in einem Punkt schneiden.

\textsc{Puissant} gibt in seinem Recueil des propositions de Géométrie einen eleganten analytischen Beweis und fügt einen geometrischen hinzu, der nicht denselben Wert hat.

\par\medskip\Needspace{5\baselineskip}
\pdfbookmark[2]{Erste handschriftliche Bemerkung.}{carnot-handschriftlich-i}
\begin{center}[\textit{Erste handschriftliche Bemerkung.}]\end{center}
\hspace*{\parindent}Sind \(\alpha,\beta,\gamma\) die komplexen Zahlen, die in natürlicher Reihenfolge zu den drei Eckpunkten eines Dreiecks gehören, \(A,B,C\) die drei Winkel und \(u\) die komplexe Zahl für den Mittelpunkt des [umschriebenen] Kreises, so gilt
\begingroup\footnotesize
\[
\begin{aligned}
2u&=\alpha+\beta+(\beta-\alpha)\operatorname{cotg}C.i=\alpha(1-i\operatorname{cotg}C)+\beta(1+i\operatorname{cotg}C)\\
&\hphantom{=\alpha+\beta+(\beta-\alpha)\operatorname{cotg}C.i}=\beta(1-i\operatorname{cotg}A)+\gamma(1+i\operatorname{cotg}A)\\
&\hphantom{=\alpha+\beta+(\beta-\alpha)\operatorname{cotg}C.i}=\gamma(1-i\operatorname{cotg}B)+\alpha(1+i\operatorname{cotg}B)
\end{aligned}
\]
\endgroup
